\section{Retrieval and Quality Control}
\label{ch:retrieval}

\subsection{Retrieval smoke test}
Each held-out query's saved CLIP image vector searches the 10,000-image training index. The evaluator sampled 20 held-out validation/test images, requested five neighbors, and excluded any same-ID result.

The test returned results for all 20 queries, with zero self-matches and mean top-1 cosine similarity 0.8260. This is a functional check, not a relevance benchmark: without judged neighbors, the similarity is not accuracy, caption quality, or evidence of visual grounding.

\begin{table}[htbp]
\centering
\caption{Recorded held-out image retrieval smoke test.}
\label{tab:retrieval-results}
\begin{tabularx}{0.85\textwidth}{@{}Xr@{}}
\toprule
Measure & Recorded result\\
\midrule
Training gallery & 10,000 images\\
Held-out queries & 20\\
Neighbors requested & 5 per query\\
Queries with results & 20 / 20\\
Self-matches & 0\\
Mean top-1 cosine similarity & 0.8260\\
\bottomrule
\end{tabularx}
\end{table}

\subsection{Quality-control results}
The pipeline QC report marks implemented checks PASS. It covers split sizes and separation, five captions per image, image decoding, exact and perceptual duplicate screening, embeddings and mappings, index readback, caption-ID exclusion, and retrieval. Perceptual screening found nine within-split candidates and zero cross-split candidates at Hamming distance 6; candidates are not proof of duplicates.

\small
\begin{longtable}{@{}p{0.30\textwidth}p{0.61\textwidth}@{}}
\caption{Selected results from the saved final pipeline quality-control report.}
\label{tab:qc-summary}\\
\toprule
\textbf{Check} & \textbf{Recorded result}\\
\midrule
\endfirsthead
\caption[]{Selected quality-control results (continued).}\\
\toprule
\textbf{Check} & \textbf{Recorded result}\\
\midrule
\endhead
\rowcolors{2}{ProjectGrey}{white}
Split sizes & Train 10,000; validation 100; test 1,000.\\
Reference captions & 5 per selected image; 55,500 records total.\\
Processed image integrity & 11,100 decoded; no missing processed images.\\
Split separation & No image-ID overlap between any pair of splits.\\
Exact duplicates & 0 duplicate groups by SHA-256.\\
Perceptual-hash screen & 9 within-split candidates; 0 cross-split candidates at Hamming distance 6.\\
Image embeddings & Expected split counts; 512 dimensions; finite and L2-normalized; manifest order checked.\\
Caption embeddings & 50,000 train vectors; 512 dimensions; unique IDs and source order checked.\\
Caption-index leakage & 0 validation/test caption IDs in the training index.\\
FAISS indexes & Image index: 10,000 vectors; caption index: 50,000; index readback passed.\\
Image retrieval & 20 of 20 queries returned results; top-5; 0 self-matches.\\
\bottomrule
\end{longtable}

\smallskip
\textit{Scope note:} The final quality-control report records zero failed, warning, and not-run checks among the checks it implements. The separate requirement-level compliance audit remains PARTIAL.


This PASS applies only to implemented pipeline checks. The separate requirement-level audit rates overall compliance PARTIAL because evidence is incomplete, including unrecorded historical preprocessing counts, partial exact caption mapping verification, and the missing unified environment snapshot. The 1,000-image test size matches the detailed assignment requirements.
